This section evaluates the hybrid through an ablation isolating each component,
the main accuracy results on the 17 SIT test cells, conformal calibration of its
uncertainty intervals, and a cross-dataset validation of the chemistry alignment
principle.

\subsection{Ablation Study and Main Results}\label{sec:solution}

We first isolate the contribution of each component, then report the full
hybrid's accuracy and calibration across the 17 test cells (the three calibration cells of Section~\ref{sec:eval} are excluded from every reported metric).

\subsubsection{Ablation Study}

\begin{table}[htbp]
\caption{Ablation study: contribution of each hybrid component (50\% split).}
\label{tab:ablation}
\resizebox{\columnwidth}{!}{%
\begin{tabular}{lrrr}
\toprule
Configuration & RMSE & Improvement vs scratch & PI Coverage \\
\midrule
Scratch LSTM             & $0.043 \pm 0.050$ & (baseline)  & n/a   \\
Zero-shot (MIT + FP, no FT) & $0.025 \pm 0.026$ & 38.2\%$^\ddag$   & 77.3\%\\
Pretrain + FT (no fing.) & $0.017 \pm 0.016$ & 47.7\%$^\ddag$   & 79.8\%\\
Full hybrid (proposed)   & $\mathbf{0.014 \pm 0.013}$ & \textbf{53.5\%}$^\ddag$  & 82.1\%\\
\bottomrule
\multicolumn{4}{l}{$^\ddag$ Wilcoxon $p<0.001$ vs scratch.}
\end{tabular}%
}
\end{table}

\begin{figure}[htbp]
\centering
\includegraphics[width=\columnwidth]{fig6_ablation}
\caption{Ablation study: mean RMSE per configuration (50\% split).}
\label{fig:ablation}
\end{figure}

Table~\ref{tab:ablation} and Figure~\ref{fig:ablation}, which plots the same
mean RMSE per configuration, show that every component contributes and the
effects compose monotonically: zero-shot transfer (pretrained model with the
fingerprint-derived initial state, \emph{no} fine tuning on target data)
already achieves a 38.2\% improvement over scratch and beats scratch on all
17 test cells; adding fine tuning without the fingerprint reaches 47.7\%;
and the full hybrid reaches 53.5\% (all Wilcoxon $p<0.001$).
Three observations stand out.
First, the zero-shot configuration is the capability no other evaluated
method offers: RMSE 0.025 for a cell with \emph{no} target-cell training,
from 30 observed cycles, recovering roughly 70\% of the full hybrid's
benefit before any adaptation is possible.
Second, the fingerprint and fine tuning are complementary rather than
redundant: the fingerprint on top of pretraining+fine tuning still reduces
RMSE from 0.017 to 0.014.
Third, a shuffled-fingerprint control clarifies the mechanism: replacing each
test cell's fingerprint with another cell's raises full-hybrid RMSE only from
0.0141 to 0.0146.
The conditioning therefore works chiefly by placing the LSTM in the
\emph{pretrained population's} operating regime (any in-distribution
early-life shape suffices) rather than by encoding a precise cell identity,
which is consistent with the two mechanisms remaining complementary
(Section~\ref{sec:mechanism}).
Prediction-interval coverage improves in the same monotone order
(77.3\%~$\rightarrow$~79.8\%~$\rightarrow$~82.1\%), so with a
well-selected pretrained checkpoint there is no accuracy versus calibration
trade-off among the configurations.

\subsubsection{Hybrid Accuracy Across Operating Conditions}

\begin{table}[htbp]
\caption{Full hybrid results by operating-condition group (50\% split).}
\label{tab:hybrid_main}
\resizebox{\columnwidth}{!}{%
\begin{tabular}{lrrrrr}
\toprule
Group & $n$ & Hybrid RMSE & Scratch RMSE & Improvement & PI Coverage \\
\midrule
G1 (25\,\textcelsius, 1C) & 9 & $0.015\pm0.012$ & $0.035\pm0.034$ & 44.0\% & 80.3\% \\
G2 (40\,\textcelsius, 1C) &  3 & $0.005\pm0.002$ & $0.016\pm0.013$ & 56.8\% & 79.0\% \\
G3 (40\,\textcelsius, 2C) &  5 & $0.019\pm0.016$ & $0.073\pm0.077$ & 68.8\% & 87.1\% \\
\midrule
All SIT (test)            & 17 & $0.014\pm0.013$ & $0.043\pm0.050$ & 53.5\% & 82.1\% \\
\bottomrule
\end{tabular}%
}
\end{table}

The full hybrid model achieves a statistically significant improvement over
scratch LSTM across the 17 test cells (Wilcoxon $W=0$, $p=1.5\times10^{-5}$;
hybrid better in 17/17 cells).
Table~\ref{tab:hybrid_main} breaks the hybrid and scratch RMSE, the mean
improvement, and the PI coverage down by operating-condition group.
Figure~\ref{fig:hybrid_percell} shows the per-cell results: the top panel
compares each cell's scratch and hybrid RMSE, while the bottom panel gives the
corresponding per-cell improvement of the hybrid over scratch (positive for
every test cell; the smallest improvement is 14.9\% on cell 101-1).

On the one-step metric itself, every method is bunched on a saturated scale:
the hybrid's mean RMSE is 0.0212, 0.0141 and 0.0156 at the 30\%, 50\% and 70\%
splits (Figure~\ref{fig:method_comparison}), the per-cell-tuned DE-LSTM reaches
0.0118 and 0.0105 at 50\% and 70\%, and even a trivial linear autoregressor sits
near 0.005, differences that carry little information on this saturated benchmark
(Section~\ref{sec:baseline}). The hybrid's decisive advantage is therefore not a
marginal one-step number but a set of capabilities that no other method on the
leaderboard provides. It is the only method that operates with no target training
at all (zero-shot, 0.025), the only method that quantifies uncertainty
(Section~\ref{sec:conformal}), and the only one validated across three
manufacturers (MIT, SIT and LISHEN), two chemistries (LFP and LCO) and three
operating conditions (Section~\ref{sec:cross_dataset}), whereas PA-LSTM, GA-LSTM
and DE-LSTM provide point estimates only, for cells that already have an
established history. On the cells that matter most for deployment, the
catastrophic degraders, this advantage is also an accuracy advantage, as the
next paragraph shows.

Within G1, the hybrid is particularly effective for catastrophic degraders
(Figure~\ref{fig:prediction}):
cell 001-8 (SOH$_\text{end}=0.067$) improves 58.2\%
(0.041 vs.\ 0.098 scratch); cell 101-3 (SOH$_\text{end}=0.118$) improves
81.1\% (0.016 vs.\ 0.083).
Both gains are driven by the MIT pretraining, which has encountered
comparable extreme-degradation trajectories in the training population,
a trajectory the scratch model cannot learn from 50\% of a single
catastrophic cell's history.
Figure~\ref{fig:prediction} plots the hybrid's predicted versus actual SOH with
95\% MC-Dropout uncertainty bands on the held-out second half (50\% split) for
three representative cells, a robust G1 cell, the catastrophic degrader 001-8,
and a condition-shifted G3 cell, with the training portion shaded and the
train/test boundary dashed.
Across all three the hybrid tracks both gentle and catastrophic trajectories, and
the prediction interval widens where degradation is steepest.

For G3 test cells (40\,\textcelsius/2C), the hybrid achieves mean RMSE of
0.019, a 68.8\% improvement over the scratch LSTM given the same 50\% G3 fine
tuning data (0.073; Table~\ref{tab:hybrid_main}), by combining MIT-learned
generic LFP priors with that fine tuning.
Cold cross-condition transfer (a G1-trained model applied to G3 with no G3
data) is not a stable reference point, since Section~\ref{sec:brittleness}
showed such estimates are dominated by training-run variance; the comparison
against the equally-fine-tuned scratch model above is the fair like-for-like
measure of the hybrid's contribution.
The G3 conclusions are not an artefact of cell 003-5 (the SOH$>$1
capacity-recovery cell): excluding it, the four remaining G3 test cells give a
hybrid RMSE of 0.014 (scratch 0.039), a 65.7\% mean per-cell improvement versus
68.8\% with 003-5 included, so the artefact slightly inflates but does not
drive the result.

\begin{figure*}[htbp]
\centering
\includegraphics[width=\textwidth]{fig7_hybrid}
\caption{Hybrid vs.\ scratch LSTM on the SIT cells (50\% split).}
\label{fig:hybrid_percell}
\end{figure*}

\begin{figure*}[htbp]
\centering
\includegraphics[width=\textwidth]{fig11_prediction}
\caption{Hybrid predicted versus actual SOH with 95\% MC-Dropout uncertainty bands (50\% split).}
\label{fig:prediction}
\end{figure*}

\subsubsection{Conformal Calibration of the Uncertainty
Intervals}\label{sec:conformal}

The raw MC-Dropout intervals are approximate (Section~\ref{sec:hybrid_arch})
and under-cover in practice: pooled over the 17 test cells, the empirical
coverage of nominal 50/80/90/95/99\% intervals is 41/68/78/84/91\%
(Figure~\ref{fig:calibration}, left).
We correct this with split-conformal prediction~\cite{angelopoulos2021gentle}
using the three calibration cells already reserved for checkpoint selection:
the normalised scores $|y_t-\hat{y}_t|/\hat{\sigma}_t$ are collected on the
calibration cells, their 95th percentile $q_{95}$ replaces the Gaussian 1.96,
and test-cell intervals become $\hat{y}_t \pm q_{95}\,\hat{\sigma}_t$.
This raises 95\% coverage from 84.4\% to \textbf{93.9\%} on the test cells
with no target-cell labels beyond the calibration set
(Figure~\ref{fig:calibration}, right).
Coverage reaches at least 92\% on 15 of 17 cells; the two exceptions are the
catastrophic degrader 001-8 (67\%) and the recovery artefact 003-5 (56\%),
whose error distributions differ structurally from anything in the
calibration set, a reminder that conformal guarantees assume
exchangeability, which extreme manufacturing outliers violate.
The full hybrid is the only method in Table~\ref{tab:method_comparison} that
produces uncertainty intervals at all; with conformal correction those
intervals become trustworthy enough for BMS threshold logic on typical cells,
while flagged outliers still require conservative handling.

\begin{figure}[htbp]
\centering
\includegraphics[width=\columnwidth]{fig13_calibration}
\caption{Uncertainty calibration on the 17 test cells.}
\label{fig:calibration}
\end{figure}

\subsubsection{Uncertainty Calibration: The Chemistry Alignment Principle}

\begin{figure*}[htbp]
\centering
\includegraphics[width=\textwidth]{fig8_coverage}
\caption{MC-Dropout 95\% PI empirical coverage by dataset and group.}
\label{fig:coverage}
\end{figure*}

Figure~\ref{fig:coverage} plots the empirical 95\% PI coverage per cell (left)
and averaged per dataset group (right), and reveals a striking pattern: PI
coverage on SIT LFP cells (82.1\% mean) is nearly double that on NASA LCO
cells (43.1\%), despite using the same MC-Dropout rate (0.2) and the same
nominal 95\% target.
The most prominent structural difference is chemistry: the MIT pretraining data
are LFP cells (matched to SIT), whereas NASA cells are LCO (LiCoO$_2$~\cite{saha2007battery}; mismatched). The
datasets also differ in cell format, laboratory and cycling protocol, so
chemistry is the most likely, but not the only possible, contributor to the
coverage gap; the cross-dataset study below helps disentangle it.

This \textbf{chemistry alignment principle} has a practical implication for
BMS deployment: transfer learning preserves uncertainty calibration only
when the pretraining and deployment chemistries are matched.
Deploying a model pretrained on LFP data to LCO cells not only degrades
point estimate accuracy but also destroys uncertainty calibration, producing
overconfident predictions that could mask true risk.
Post-hoc calibration (conformal prediction~\cite{angelopoulos2021gentle}
or isotonic regression) is required when chemistry is mismatched.

\subsection{Cross-Dataset Validation of the Chemistry Alignment
  Principle}\label{sec:cross_dataset}

The analysis in Section~\ref{sec:solution} established the chemistry alignment
principle from a single source-target pair (MIT$\rightarrow$SIT).
To verify that this finding generalises, we evaluate the hybrid framework
across ten pairwise transfer scenarios spanning three LFP datasets
(MIT Severson~\cite{severson2019data}, SIT, and LISHEN 40\,Ah LFP) and one
LCO dataset (NASA 18650).
The LISHEN dataset~\cite{lishen_zenodo} provides three industrial-scale LFP
cells (40\,Ah) cycled at 2C with 20\% depth of discharge (DOD) regular cycles and full-discharge
reference tests every tenth cycle, offering a third independent LFP
pretraining source distinct from both MIT and SIT in cell format, capacity,
and cycling protocol.

\begin{table*}[htbp]
\centering
\caption{Cross-dataset transfer matrix: RMSE improvement over scratch LSTM.}
\label{tab:cross_dataset}
\begin{tabular*}{\textwidth}{@{\extracolsep{\fill}}clllrrr}
\toprule
Sc. & Source & Target & Chem. & ZS+FP & PT+FT & Hybrid \\
\midrule
\multicolumn{7}{@{}l}{\textit{LFP$\leftrightarrow$LFP (chemistry matched)}} \\
A & MIT (40\,cells) & SIT        & LFP & $+$38.2\% & $+$47.7\% & $+$53.5\% \\
G & MIT (40\,cells) & LISHEN     & LFP & $+$80.2\% & $+$66.2\% & $+$64.5\% \\
B & SIT (20\,cells) & MIT        & LFP & $+$19.6\% & $+$57.1\% & $+$59.5\% \\
H & SIT (20\,cells) & LISHEN     & LFP & $-$22.6\% & $+$46.1\% & $+$38.6\% \\
N & LISHEN (3\,cells) & MIT      & LFP & $+$14.6\% & $+$4.7\% & $+$4.1\% \\
M & LISHEN (3\,cells) & SIT      & LFP & $-$24.6\% &  $+$5.8\% &  $+$5.2\% \\
\midrule
\multicolumn{7}{@{}l}{\textit{LCO$\leftrightarrow$LFP (chemistry-mismatched)}} \\
C & MIT (LFP) & NASA (LCO) & cross & $+$15.6\% & $+$8.8\% & $+$8.9\% \\
D & SIT (LFP) & NASA (LCO) & cross & $+$4.8\% & $+$21.6\% & $+$21.6\% \\
E & NASA (LCO) & SIT (LFP) & cross & $+$26.8\% & $+$44.1\% & $+$43.1\% \\
F & NASA (LCO) & MIT (LFP) & cross & $+$37.2\% & $+$59.8\% & $+$59.5\% \\
\bottomrule
\end{tabular*}
\end{table*}

\begin{figure*}[htbp]
\centering
\includegraphics[width=\textwidth]{fig9_cross_dataset}
\caption{Cross-dataset RMSE improvement over scratch LSTM across all ten
transfer scenarios.}
\label{fig:cross_dataset}
\end{figure*}

Scenario~A (MIT$\rightarrow$SIT) repeats the main ablation of
Table~\ref{tab:ablation}, included here for reference.
In Table~\ref{tab:cross_dataset}, the ten scenarios are grouped by chemistry
alignment (LFP$\leftrightarrow$LFP matched, LCO$\leftrightarrow$LFP mismatched)
and negative values denote degradation relative to scratch.
Together with Figure~\ref{fig:cross_dataset}, the results support three findings.

\subsubsection{Fine tuning is the dependable adaptation mechanism}
Pretraining plus fine tuning improves over scratch in \emph{all ten}
scenarios ($+$4.7\% to $+$66.2\%, mean 36.2\%), and the full hybrid tracks
it closely (mean 35.9\%), being best or within one point of best in six of
ten (A, B, D, E, F and~M).
Whatever the source corpus, chemistry, or cell format, fine tuning a
pretrained model on the target cell's history is never worse than training
from scratch in these experiments, which makes it the mechanism a
practitioner can rely on when facing an unfamiliar dataset pair.
The weakest transfers (M and~N, both below 6\%) share one cause: the 3-cell
LISHEN corpus is too narrow a pretraining source for either mechanism.

\subsubsection{Zero-shot conditioning is a distinct capability, not a
universal accuracy win}
Zero-shot performance varies far more across scenarios ($-$24.6\% to
$+$80.2\%) than fine tuning does.
It is strongly positive precisely where the pretrained population resembles
the deployment population: the paper's primary scenario~A
(MIT$\rightarrow$SIT, $+$38.2\%) and scenario~G (MIT$\rightarrow$LISHEN,
$+$80.2\%), both large LFP corpora transferring to LFP targets.
It fails when the source poorly spans the target (H and M, both negative:
SIT's 20 heterogeneous cells and LISHEN's 3 cells do not cover each other's
regimes), because the fingerprint places a cell within the
\emph{pretraining} population's degradation distribution, and that placement
misleads when the distribution does not contain the target's behaviour.
The correct reading is therefore about regime, not rank: zero-shot
conditioning is the only option during the first weeks of deployment, before
any fine-tuning history exists, and it is trustworthy exactly when the
practitioner can match pretraining chemistry and corpus breadth to the
deployment population, precisely the condition a battery manufacturer
qualifying one product line can guarantee.

\subsubsection{Chemistry mismatch reduces accuracy gains and destroys
calibration}
Across the four cross-chemistry scenarios (C to F), hybrid improvement ranges
from 8.9\% to 59.5\% (mean 33.3\%), against 38.6\% to 64.5\% (mean
54.0\%) for the adequately resourced matched-LFP scenarios (A, B, G, H).
The sharper signal is calibration: the PI coverage analysis from
Section~\ref{sec:solution} (82\% LFP$\rightarrow$LFP vs 43\%
LFP$\rightarrow$LCO) shows that the RMSE improvement metric understates the
damage from chemistry mismatch.
The model can remain reasonably accurate in point estimates while becoming
severely overconfident in its uncertainty intervals, the failure mode most
dangerous for BMS deployment, and one invisible to every uncertainty-free
method in Table~\ref{tab:method_comparison}.
